\section{Piattaforma Numerica}

% 
\begin{frame}
\frametitle{Piattaforma Numerica}

\begin{center}
 \vspace{0.3cm}
 %     \begin{center}
%      
%    \large La piattaforma numerica  semplificato e versatile}
%    che permette l'integrazione e l'accoppiamento dei codici computazionali di \textbf{\textcolor{bleudefrance}{neutronica}} e
%    di \textbf{\textcolor{mygreen}{termoidraulica}}.
%    
%     \end{center}
 La piattaforma numerica consiste in un \textbf{ambiente computazionale}, basato sul progetto \textit{Salom\'e}, e permette la  
%  coordinare i codici computazionali
 \textbf{gestione} e \textbf{condivisione} delle strutture dati 
 tra i moduli di \textbf{\textcolor{bleudefrance}{neutronica}} e di \textbf{\textcolor{mygreen}{termoidraulica}} \\
 \end{center}
%  \vspace{0.5cm}
%  \begin{center}
  \begin{columns}\hspace{0.5cm}
 \begin{column}{0.42\textwidth}
 \begin{block}{Obiettivo}
%  \begin{center}
   \begin{itemize}
    \item Consentire l'\textbf{accoppiamento} tra la 
     neutronica
     e la termoidraulica
    \item Ottenere la soluzione di un \textbf{problema multifisico} in un ambiente semplificato e versatile
   \end{itemize}
  \end{block}
%   \begin{block}{Scopo dell'accoppiamento neutronico-termoidraulico}
%   \begin{itemize}
%   \item Scambiare la potenza termica ottenuta dal modulo di neutronica con il
%   modulo di termoidraulica
%   \item A sua volta, il modulo di termoidraulica restituisce la distribuzione del campo di temperatura
%   al modulo di neutronico al fine di correggere il flusso neutronico a seguito delle variazioni di tempertura
%   \end{itemize}
%   \end{block}
\end{column}
\begin{column}{0.5\textwidth}
\begin{minipage}{\textwidth}
\begin{tikzpicture}[overlay]
 \node[xshift=25mm,yshift=32mm, text width = 4cm](sigma) {\textbf{\textcolor{mygreen}{$T_f,T_c$}}};
 \node[xshift=45mm,yshift=-2mm, text width = 4cm](sigma) {\textbf{\textcolor{bleudefrance}{$\phi$, $P_{th}$}}};
 \node[xshift=63mm,yshift=30mm, text width = 4cm](TF) {\textbf{\textcolor{red}{$\Sigma_f,\Sigma_a,\Sigma_s$}}};
\end{tikzpicture}
\scalebox{.7}{\smartdiagramset{set color list={red, bleudefrance, mygreen}, }\smartdiagram[circular diagram:clockwise]{Database Nucleare (Dragon), Simulazione Neutronica (Donjon), Analisi Termo Idraulica (FEMuS)}}
\end{minipage}
\end{column}
\end{columns}
\vspace{0.2cm}
\begin{center}
 Entrambi i codici ottengono la soluzione numerica attraverso \\ il \textbf{metodo degli elementi finiti}
\end{center}


\end{frame}
% %%%%%%%%%%%%%%%%%%%%%%%%%%%%%%%%%%%%%%%%%%%%%%%%%%%%%%%%%%
% \begin{frame}
% Equazione della diffusione neutronica
% \begin{multline*}
%  -\nabla \cdot(D_g(\vec{r})\nabla \phi_g(\vec{r}))+\Sigma_g(\vec{r})\phi_g(\vec{r})=\\
%  =\sum_{h=1}^G \Sigma_{h \rightarrow g}(\vec{r}) 
% \phi_h(\vec{r})+\frac{\chi_g(\vec{r})}{k_{eff}} \sum_{h=1}^G \nu 
% \Sigma_{f,h}(\vec{r})\phi_h(\vec{r})   \,.
% \end{multline*}
% Equazione del calore
% \begin{equation*}
%  \rho c_v\frac{dT}{dt}=\nabla \cdot (\kappa \nabla T) + \mu \Phi + \rho 
% \vec{g}\cdot \vec{\upsilon} + \dot{Q} \,.
% \end{equation*}
% \end{frame}
% %%%%%%%%%%%%%%%%%%%%%%%%%%%%%%%%%%%%%%%%%%%%%%%%%%%%%%%%%%5
\subsection{Codice di Neutronica}

\begin{frame}
 \frametitle{Modulo di Neutronica}
 
 \begin{center} \vspace{0.2cm} 
 \begin{block}{}
   \begin{itemize}
     \item \textbf{\textcolor{red}{Dragon}}: Codice di cella risolve le \textbf{equazioni del trasporto multigruppo} 
     allo scopo
     di ottenere le \textit{sezioni d'urto macroscopiche $\Sigma_i$}  e di generare il \textbf{database nucleare}
   \end{itemize}
   \end{block}
   \begin{block}{}
   \begin{itemize}
    \item \textbf{\textcolor{bleudefrance}{Donjon}}: Codice di reattore che simula il comportamento dei neutroni sull'intero
          reattore, ottenendo la distribuzione di \textbf{flusso neutronico} $\phi_g$ e di \textbf{potenza termica}. 
%           attraverso la risoluzione
% 	  dell'equazione di diffusione o di equazioni $SP_N$.
      \end{itemize}	  
  \end{block}

   \underline{Equazione multigruppo della diffusione neutronica}
\begin{multline*}
 -\nabla \cdot(D_g(\vec{r})\nabla \phi_g(\vec{r}))+\Sigma_g(\vec{r})\phi_g(\vec{r})=\\
 =\sum_{h=1}^G \Sigma_{h \rightarrow g}(\vec{r}) 
\phi_h(\vec{r})+\frac{\chi_g(\vec{r})}{k_{eff}} \sum_{h=1}^G \nu 
\Sigma_{f,h}(\vec{r})\phi_h(\vec{r})   \,.
\end{multline*}
  \end{center}

  \end{frame}
 
 \subsection{Codice di Termoidraulica}
 
 \begin{frame}
  \frametitle{Modulo di Termoidraulica}
   \begin{center}
%    \vspace{0.5cm}
   \begin{block}{}
   \begin{itemize}
    \item \textbf{\textcolor{mygreen}{FEMuS}}: Codice CFD che risolve le equazioni di termofluidodinamica
    con possibilità di aggiungere modelli di turbolenza.   
    \end{itemize}
      \end{block}
      
    \underline{Equazione del calore per fluidi incomprimibili}
    \begin{equation*}
 \rho c \frac{dT}{dt}=\nabla \cdot (\kappa \nabla T) + \mu \Phi + \rho 
\vec{g}\cdot \vec{\upsilon} + \dot{Q}_{neutronica} 
\end{equation*}



 \underline{Proprietà termodinamiche del piombo}
 \begin{table}[!htb]
\centering
 \begin{tabular}{lccc}
\hline
Parametro & Simbolo & Unità & Valore \\
\hline
Densità & $\rho$ & kg/m$^3$ & 11340 \\
Conducibilità termica & $\kappa$ & W/(m$\cdot$K) & 35 \\
Calore specifico & $c$ & J/(kg$\cdot$K) & 129 \\
Viscosità dinamica & $\mu$ & Pa$\cdot$s & 0.001844 \\
\hline
\end{tabular}
\end{table}

  \begin{itemize}
 \item Non è necessario risolvere Navier-Stokes (Canali chiusi)
 \item Conducibilità termica mediata su tutti i materiali dell'assembly 
\end{itemize}
%   \item \begin{block}{FEMuS}
%          Codice di termoidraulica che risolve le principali equazioni di termodinamica.
%         \end{block}
%   \end{itemize}
%   In supporto alla neutronica, il

 \end{center}

\end{frame}

%%%%%%%%%%%%%%%%%%%%%%%%%%%%%%%%%%%%%%%%%%%%%%%%%%%%%%%%%%%%%%%%%%%%

\begin{frame}
\frametitle{Logica dell'accoppiamento neutronico-termoidraulico}
\vspace{0.1cm}
\begin{columns}\hspace{0.5cm}
 \begin{column}{0.35\textwidth}
 
% 
\begin{center} 
  \begin{itemize}
   \item Librerie dati nucleari \textbf{\textcolor{orange}{JEFF-3.1.2-SHEM315}}
   \item Codice di cella\\ \textbf{\textcolor{red}{Dragon}}
   \item Codice di reattore\\ \textbf{\textcolor{bleudefrance}{Donjon}}
   \item Codice di termoidraulica \textbf{\textcolor{mygreen}{FEMuS}}
  \end{itemize}
  
\end{center}
%  \end{block}
 \end{column}

 \begin{column}{0.72\textwidth}
  \begin{center}
 \includegraphics[scale=0.55]{../../thesis/figure/flowchart/dataflow_platform}
 \end{center}
 \end{column}
 \end{columns}

\end{frame}

% % % % %  %%%%%%%%%%%%%%%%%%%%%%%%%%%%%%%%%%%%%%%%%%%%%%%%%%%%%%%%%%%%%%%%%%%%%%
% % % % %  \section{Modellazione del reattore}
% % % % %  \begin{frame}
% % % % %  \frametitle{Creazione del database nucleare}
% % % % %  \vspace{0.5cm}
% % % % %  
% % % % %  \begin{center}
% % % % %   Il \textbf{database nucleare} contiene le informazioni sulle sezioni d'urto macroscopiche $\Sigma_T$,
% % % % %   omogeneizzate su ciascuna cella e condensate a 33 gruppi di energia, degli elementi che compongono il
% % % % %   reattore.
% % % % %  \end{center}
% % % % %   \begin{columns}
% % % % %   
% % % % %     \begin{column}{0.6\textwidth}
% % % % % %       \begin{center}
% % % % % %     \includegraphics[scale=1]{../../thesis/figure/itex/unitcells_wq}
% % % % % %     
% % % % % %    \end{center}
% % % % %      Parametri per le interpolazione
% % % % %      \begin{itemize}
% % % % %       \item Temperatura del combustibile $T_f$
% % % % %       \item Temperatura del refrigerante $T_c$
% % % % %       \item Tempo relativo al bruciamento $t_{burnup}$
% % % % %      \end{itemize}
% % % % % 
% % % % % 
% % % % %    
% % % % %   \end{column}
% % % % %   
% % % % % 
% % % % %   
% % % % %   \begin{column}{0.4\textwidth}
% % % % %    \begin{center}
% % % % %     \includegraphics[scale=0.26]{../../thesis/figure/itex/fassembly}\\
% % % % %     Assembly di combustibile
% % % % %    \end{center}
% % % % %    \end{column}
% % % % %    
% % % % %   
% % % % %   
% % % % %  \end{columns}
% % % % %  \begin{block}{}
% % % % %    \begin{center}
% % % % %   La sua realizzazione permetterà al codice di reattore \textcolor{bleudefrance}{\textbf{Donjon}} di correggere le sezioni d'urto a seguito delle variazioni
% % % % %   di temperatura.
% % % % %  \end{center}
% % % % % 
% % % % %  \end{block}
% % % % %  \end{frame}
% % % % %  
%%%%%%%%%%%%%%%%%%%%%%%%%%%%%%%%%%%%%%%%%%%%%%%%%%%%%%%%%%%%%%%%%%%%%%%%%%%%%%%%%%%%%%%%%%%%%%%%%%%%%%%%%%%%% 
% % % %  %%%%%%%%%%%%%%%%%%%%%%%%%%%%%%%%%%%%%%%%%%%%%%%%%%%%%%%%%%%
% % % %  \begin{frame}
% % % %  
% % % %     \begin{columns}
% % % %   
% % % %     \begin{column}{0.6\textwidth}
% % % %  \small{
% % % %  \begin{tabular}{lcc}
% % % % \hline
% % % % Parametro & Unità & Valore \\
% % % % \hline
% % % % N$^\circ$ totale di esagoni & - & 169 \\
% % % % N$^\circ$ barre di combustibile & - & 126 \\
% % % % N$^\circ$ elementi di acciaio & - & 43 \\
% % % % Lato di un esagono & mm & 8.00208 \\
% % % % Apotema di un esagono & mm & 6.93 \\
% % % % Distanza tra barre di combustibile & mm & 13.86 \\
% % % % Raggio barra di combustibile & mm & 5.25 \\
% % % % \hline
% % % % \end{tabular}
% % % %    }
% % % %   \end{column}
% % % %   
% % % % 
% % % %   
% % % %   \begin{column}{0.4\textwidth}
% % % %   
% % % %    \begin{center}
% % % %     \underline{Esagono di combustibile}
% % % %     \includegraphics[scale=1.2]{../../thesis/figure/itex/unitcells_wq}\\
% % % % 
% % % %    \end{center}
% % % %    \end{column}
% % % %    
% % % %   
% % % %   
% % % %  \end{columns}
% % % %   
% % % %   
% % % %  \end{frame}
% % % % 
% % % %  %%%%%%%%%%%%%%%%%%%%%%%%%%%%%%%%%%%%%%%%%%%%%%%%%%%%%%%%%%%
 \begin{frame}
 \frametitle{Geometria complessiva del reattore}
 \vspace{0.3cm}
 \begin{columns}
  
  \begin{column}{0.48\textwidth}
%   \begin{itemize}
%    \item[-] \textcolor{yellow}{$\circ$} Combustibile MOX arricchito al \underline{20.5$\%$}
%    \item[-] \textcolor{red}{$\circ$} Combustibile MOX arricchito al \underline{26.2$\%$}
%    \item[-] \textcolor{green}{$\circ$} Barre di controllo
%    \item[-] \textcolor{purple}{$\circ$} Barre di sicurezza
%   \end{itemize}
   \begin{center}
   
  \underline{Sezione radiale del nocciolo}
    \includegraphics[scale=0.3]{../../thesis/figure/itex/hexagon}
   \end{center}

  \end{column}
  
    \begin{column}{0.48\textwidth}
     \begin{center}
     \underline{Sezione assiale del reattore}
    \includegraphics[scale=0.28]{../../thesis/figure/itex/rectangle_wq}
   \end{center}
   

  \end{column}
  
 \end{columns}
 \vspace{0.3cm}
 \begin{center}
 \begin{itemize}
  \item  Arricchimento combustibile MOX interno {\underline{20.5$\%$}} (in giallo), 
    esterno {\underline{26.2$\%$} (in rosso)}
  \item 12 barre di controllo (in verde) e 4 di sicurezza (in viola)
 \end{itemize}
 \end{center}

 \end{frame}
 %%%%%%%%%%%%%%%%%%%%%%%%%%%%%%%%%%%%%%%%%%%%%%%%%%%%%%%%%%%%%
 
  %%%%%%%%%%%%%%%%%%%%%%%%%%%%%%%%%%%%%%%%%%%%%%%%%%%%%%%%%%%%%
% % % %  \begin{frame}
% % % %  \frametitle{Modellazione complessiva del reattore}
% % % %  \begin{columns}
% % % %   
% % % %   \begin{column}{0.48\linewidth}
% % % %      \begin{center}
% % % %   \underline{Sezione assiale del reattore}
% % % %     \includegraphics[scale=0.35]{../../thesis/figure/itex/rectangle_wq}
% % % %    \end{center}
% % % %   \end{column}
% % % %   
% % % %     \begin{column}{0.48\linewidth}
% % % %      \begin{center}
% % % %      \underline{Reattore tridimensionale}
% % % %        \includegraphics[scale=0.3]{../../thesis/figure/result_tred/new/tred-reactor}
% % % %    \end{center}
% % % %   \end{column}
% % % %   
% % % %  \end{columns}
% % % %  



% % % % %  \end{frame}
 %%%%%%%%%%%%%%%%%%%%%%%%%%%%%%%%%%%%%%%%%%%%%%%%%%%%%%%%%%%%%
%  \begin{frame}
%  \frametitle{Codice di termoidraulica}
%  \begin{columns}
%   
%   \begin{column}{0.48\textwidth}
%      \begin{center}
%     \includegraphics[scale=0.4]{../../thesis/figure/itex/tempeq}
%    \end{center}
%   \end{column}
%   
%     \begin{column}{0.48\textwidth}
%      \begin{center}
%     \includegraphics[scale=0.35]{../../thesis/figure/itex/canale}
%    \end{center}
%   \end{column}
%   
%  \end{columns}
%  \end{frame}